\documentclass[a4paper,10pt]{article}
\usepackage{graphicx} % Required for inserting images
\usepackage[utf8]{inputenc}
\usepackage[T1]{fontenc}
\usepackage{enumitem}
\usepackage{geometry}
\geometry{left=1.5cm,right=1.5cm,top=1cm,bottom=1cm}
\usepackage{fontawesome5}
\usepackage{hyperref}
\usepackage{xcolor}
\usepackage{tabularx}
\usepackage{array}
\usepackage{setspace}
\usepackage[sfdefault,light]{roboto}  % Définit Roboto comme police sans-serif
\renewcommand{\familydefault}{\sfdefault}  % L’applique comme police par défaut

\pagestyle{empty}
\setlength{\parindent}{0pt}
\definecolor{linkcolor}{HTML}{1A73E8}

\hypersetup{
  colorlinks=true,
  urlcolor=linkcolor,
}

\begin{document}
    
    %=========== HEADER ===========
    \begin{tabularx}{\textwidth}[t]{@{}X r@{}}
        \begin{tabular}[t]{@{}l@{}}
            {\fontsize{18}{18}\selectfont\textbf{MARSEL BAKASHOV}} \\[0.4em]
            {\fontsize{12}{12}\selectfont\textit{Maîtrise en apprentissage automatique — Mila}} \\[0.3em]
            {\fontsize{11}{11}\selectfont Finance quantitative / Risque / Stage mai--août 2027}
        \end{tabular}
        &
        \begin{tabular}[t]{@{}ll@{}} 
            \begin{tabular}[t]{@{}l@{}} 
                \faGlobe\hspace{0.6em} \href{https://marselbakashov.netlify.app/}{Portfolio} \\[0.7em]
                \reflectbox{\faPhone}\hspace{0.6em} +1 (514) 664-9320 \\[0.7em]
                \faMapMarker*\hspace{0.6em} Montréal, Canada
            \end{tabular}
            &
            \begin{tabular}[t]{@{}l@{}} 
                \faEnvelope\hspace{0.6em} \href{mailto:m.bakashov11@gmail.com}{m.bakashov11@gmail.com} \\[0.7em]
                \faGithub\hspace{0.6em} \href{https://github.com/Mbaka11}{GitHub} \\[0.7em]
                \faLinkedin\hspace{0.6em} \href{https://www.linkedin.com/in/marsel-bakashov-813643213/}{LinkedIn}
            \end{tabular}
        \end{tabular}
    \end{tabularx} \\[0.3em]
    
    %===== SUMMARY & SKILLS (2 colonnes) =====
    \noindent
    \begin{tabularx}{\textwidth}[t]{@{}p{0.48\textwidth}@{\hspace{2em}}p{0.48\textwidth}@{}}
      % Bloc gauche – À PROPOS
      \parbox[t]{\linewidth}{
        \makebox[0pt][l]{\textbf{\large À PROPOS}}%
        \hspace{6em}\rule{\dimexpr\linewidth-6.2em}{0.4pt} \\[0.5em]
            Candidat à la maîtrise en apprentissage automatique à Mila, formé en génie logiciel. Expérience en tarification dynamique, optimisation sous contraintes et modélisation prédictive. Maîtrise de Python, SQL et des plateformes de données pour développer, simuler et valider des modèles quantitatifs orientés décision.
      }
      &
      % Bloc droit – COMPÉTENCES
      \parbox[t]{\linewidth}{
        \makebox[0pt][l]{\textbf{\large COMPÉTENCES}}%
        \hspace{9em}\rule{\dimexpr\linewidth-9.5em}{0.5pt} \\[1em]
        \begin{tabular}{@{}p{2.7cm}p{\dimexpr\linewidth-3.2cm}@{}}
          \textbf{Programmation:} & Python, SQL, R, C++, MATLAB \\[0.2em]
          \textbf{Modélisation:} & NumPy, Pandas, Scikit-learn, PyTorch, Monte Carlo, RL, CVaR \\[0.2em]
          \textbf{Données / BI:} & Bloomberg, Databricks, Snowflake, Power BI, Tableau, MLflow
          \\[0.2em]
          \textbf{Langues:} & Français, Anglais, Russe, Espagnol, Coréen
        \end{tabular}
      }
    \end{tabularx}
    
    %===== EDUCATION (1 colonne) =====
    \vspace{0.2em}
    \textbf{\large ÉDUCATION} \hspace{0.2em}\rule{\dimexpr\linewidth-8em}{0.4pt} \\[0em]
    
    \begin{tabularx}{\textwidth}{@{}p{3cm} X r@{}}
        1/2026 -- 2/2027 & 
        \textbf{Maîtrise pro en apprentissage automatique} 
        & \textbf{\textit{Université de Montréal / McGill  - MILA}} \\ 
        & GPA: 4.0 / 4.0 & 
        Montréal, Canada \\[0.2em]
        
        9/2021 -- 12/2025 & 
        \textbf{Baccalauréat en génie logiciel} 
        & \textbf{\textit{Polytechnique Montréal}} \\ 
        & GPA: 3.4 / 4.0 & 
        Montréal, Canada \\[0.2em]
        
        2/2025 -- 6/2025 & 
        \textbf{Baccalauréat en génie informatique} 
        & \textbf{\textit{Korea University College of Engineering}} \\ 
        & Échange académique & 
        Séoul, Corée du Sud
    \end{tabularx}
    
    %===== EXPÉRIENCE =====
    \vspace{0.2em}
    \textbf{\large EXPÉRIENCE} \hspace{1em}\rule{\dimexpr\linewidth-8.2em}{0.4pt} \\[0em]
    
    \begin{tabularx}{\textwidth}{@{}p{3cm} X@{}}
        9/2026 -- 12/2026 & 
        \textbf{Stagiaire en science des données et ML} \hfill \textbf{\textit{McKesson}}\\[0.2em]
        & \textbullet\ \hangindent=0.5em Conçu des agents IA et automatisé des flux d’analyse de données pour produire des synthèses exécutives et appuyer la prise de décision. \\
        & \texttt{Power BI / Power Automate / Python / Azure / Claude Code / Snowflake} \\[0.3em]

        1/2026 -- 5/2026 & 
        \textbf{Stagiaire en science des données et RL} \hfill \textbf{\textit{Intact Assurance}}\\[0.2em]
        & \textbullet\ \hangindent=0.5em Conçu et validé par simulation des politiques de tarification par RL avec optimisation primal--dual et mesure du risque CVaR, sous contraintes réglementaires et d’affaires, afin d’accroître la rentabilité et de limiter la volatilité.\\
        & \texttt{Python / PyTorch / RL / CVaR / MLflow / Databricks / Snowflake} \\[0.3em]
        
        6/2025 -- 9/2025 & 
        \textbf{Stagiaire en science des données et ML} \hfill \textbf{\textit{Pratt \& Whitney}}\\[0.2em]
        & \textbullet\ \hangindent=0.5em Développé des modèles prédictifs de durée de vie restante des moteurs pour appuyer la planification de la maintenance et réduire de 15 \% les interruptions imprévues.\\
        & \texttt{Python / Scikit-learn / PyTorch / Pandas / NumPy / MLflow / Jupyter} \\[0.3em]
        
        5/2024 -- 9/2024 & 
        \textbf{Stagiaire en science des données} \hfill \textbf{\textit{Radio-Canada / CBC}}\\[0.2em]
       & \textbullet\ \hangindent=0.5em Conçu un système d’indexation centralisant 15 M+ d’éléments pour la recherche sémantique et lexicale, utilisé par plus de 100 employés. \\
        & \texttt{ELK / ETL / Airflow / Python / Azure AI Search / Hugging Face / Embeddings} \\[0.3em]
    \end{tabularx}
    
    %===== PROJETS =====
    \vspace{0.2em}
    \textbf{\large PROJETS} \hspace{1em}\rule{\dimexpr\linewidth-7em}{0.4pt} \\[0em]
    
    \begin{tabularx}{\textwidth}{@{}p{3cm} X@{}}
        Projet personnel & \textbf{Finance quantitative — Monte Carlo, GBM et produits dérivés} \hfill \textbf{\textit{\href{https://github.com/Mbaka11/Quant-Finance-Learning-Notes}{GitHub}}} \\
        & \textbullet\ \hangindent=0.5em Développé des notebooks et scripts reproductibles de simulation Monte Carlo, de GBM et de valorisation de produits dérivés, avec calibration et validation sur données de marché. \\
        & \texttt{Python / NumPy / Pandas / Matplotlib / Monte Carlo / GBM / Finance quantitative} \\[0.2em]

        Projet personnel & \textbf{GravityGuyML — Apprentissage par renforcement profond} \hfill \textbf{\textit{\href{https://github.com/Mbaka11/Gravity-Guy-ML}{GitHub}}} \\
        & \textbullet\ \hangindent=0.5em Conçu un environnement de simulation \texttt{Gymnasium} avec génération procédurale et implémenté plusieurs agents profonds (\texttt{DQN}, \texttt{PPO}, \texttt{SAC}, \texttt{A2C}). \\
        & \texttt{Python / Simulation / RL / TensorFlow / Pygame} \\[0.2em]
        
        Projet personnel & 
        \textbf{Prévision d’épidémies à partir de données climatiques} \hfill \textbf{\textit{\href{https://github.com/Mbaka11/20251R0136COSE47101}{GitHub}}}\\[0.2em]
        & \textbullet\ \hangindent=0.5em Entraîné et comparé des modèles de classification (régression logistique, SVM, arbres de décision), puis évalué leur performance avec les métriques F1 et AUC-ROC. \\
        & \texttt{Python / Pandas / Scikit-learn / Matplotlib / Jupyter / FP-Growth} \\
    \end{tabularx}

    %===== AUTRES IMPLICATIONS =====
    \vspace{0em}
    \textbf{\large IMPLICATIONS PARASCOLAIRES} \hspace{1em}\rule{\dimexpr\linewidth-20em}{0.4pt} \\[0em]
    
    \begin{tabularx}{\textwidth}{@{}p{3cm} X@{}}
        1/2024 -- 8/2024 & 
        \textbf{Assistant d’enseignement – Mathématiques discrètes} \hfill \textbf{\textit{Polytechnique Montréal}} \\
        & \textbullet\ Accompagné plus de 80 étudiants en logique et graphes lors de séances de révision.\\[0.1em]
        
        1/2024 -- 5/2024 & 
        \textbf{Directeur des commandites} \hfill \textbf{\textit{PolyAI Hackathon}} \\
        & \textbullet\ Obtenu 22~000~\$ en commandites et coordonné une équipe de partenariats. 
    \end{tabularx}

\vspace{0cm}

\end{document}